\section{Algebraic objects: from quotient rings to signature bytes}
\label{sec:algebraic-object}

\subsection{The base rings and mode~2 parameters}

In the pinned specification,
\[
  R=\ZZ[X]/(X^{256}+1),\qquad
  R_q=\ZZ_q[X]/(X^{256}+1),\qquad q=64513
\]
and the mode~2 module dimensions are \((k,\ell)=(2,4)\).
The small-secret-coefficient bound is \(\eta=1\), and the binary challenge
has Hamming weight \(\tau=58\) \cite{haetae-spec}.
Each \(R_q\) element is an equivalence class of polynomials of degree
\(<256\); choosing coefficient vectors identifies it with \(\ZZ_q^{256}\)
as a set.

The specification chooses \(q\equiv1\pmod{512}\), so \(X^{256}+1\) splits
completely over \(\FF_q\) and an NTT is available.
On paper, one can therefore consider the Chinese remainder theorem isomorphism
\[
  R_q\simeq\prod_{i=1}^{256}\FF_q
\]
The repository's top-level implementation-security interface nevertheless
retains \identifier{obl_public_key_algebra_ntt} as a separate obligation.
Existence of this isomorphism and correctness of actual NTT routines must not
be treated as the same claim.

\begin{caution}[Separate mathematical background from implementation theorems]
The ring/module description in this section explains the pinned HAETAE
specification. The implementation theorem that “actual Jasmin NTT computes
this isomorphism correctly” requires a separate proposition and premises.
This guide applies \statusproved\ only to named theorems under
\sourcepath{easycrypt/}.
\end{caution}

\subsection{Viewing KeyGen as a module map}

In the paper's KeyGen data flow, a seed generates
\(a_{\mathrm{gen}}\in R_q^k\) and
\(A_{\mathrm{gen}}\in R_q^{k\times(\ell-1)}\).
Small elements \(s_{\mathrm{gen}}\in R^{\ell-1}\) and
\(e_{\mathrm{gen}}\in R^k\) are sampled, and
\[
  b=a_{\mathrm{gen}}+A_{\mathrm{gen}}s_{\mathrm{gen}}
      +e_{\mathrm{gen}}\in R_q^k
\]
is computed. The coefficientwise decomposition
\[
  b=b_0+2b_1,\qquad
  (b_0,b_1)=
  (\operatorname{LowBits}_{vk}(b),
   \operatorname{HighBits}_{vk}(b))
\]
then compresses the verification key to \(vk=(\mathit{seedA},b_1)\).
The packed secret-key representation includes those bytes as a prefix so that
Sign can read the same verification-key components.

The first composition seam closed in this repository is this inclusion
relation, not the complete KeyGen distribution.
Let \(S\) be the set of secret-key byte arrays and \(V\) the set of
verification-key arrays, and define
\[
  \pi_{992}:S\longrightarrow \Bytes^{992},\qquad
  \pi_{992}(s)=s[0..992)
\]
For a returning execution of actual mode~2 internal KeyGen with result
\((vk,sk)\),
\[
  \pi_{992}(sk)=\pi_{992}(vk)
\]
The array on the right is larger; only its first 992 bytes are compared.

This can also be expressed precisely as a quotient. Define on \(S\)
\[
  s\sim_{992}s'\quad\Longleftrightarrow\quad
  s[0..992)=s'[0..992)
\]
Then \(\pi_{992}\) classifies \(S/{\sim_{992}}\).
The current theorem states that the \(\sim_{992}\)-equivalence class of the
packed secret key agrees with the public-key representation. It does not
state equality of the entire secret keys.

\subsection{The Week 16 KeyGen snapshot congruence}

Week~16's transparent harness
\identifier{ActualM23MatrixFinalizeSnapshot.run} calls actual
\identifier{_kp_m23_matrix} with mode~2 dimensions \((2,3)\), stores its
returned \(bp\) as \(\mathit{pre\_bp}\), and calls actual
\identifier{_keypair_finalize_m23} on 512 active words
\cite{week16-mincore-plan,week16-kg-report,week16-kg-ntt-mul-report}.
At coefficient \(i\), let
\[
  r_i=(\mathit{pre\_bp}_i+s_{2,i}+a_i)\bmod q
\]
The compiled semantic snapshot gives
\[
  r_i=2b_{1,i}+b_{0,i},\qquad -1\le b_{0,i}\le1,
  \qquad s'_{2,i}=s_{2,i}-b_{0,i}
\]
Moreover,
\[
  2(a_i-2b_{1,i})+2\mathit{pre\_bp}_i+2s'_{2,i}
  \equiv0\pmod{2q}
\]
holds at every active coefficient. Here \(\mathit{pre\_bp}\) is the return
snapshot of the actual matrix helper; it has not yet been identified with
the paper's \(A_{\mathrm{gen}}s_{\mathrm{gen}}\).

\begin{caution}[Separate the snapshot congruence from the paper's KeyGen equation]
\theoremname{actual_m23_matrix_finalize_semantic_snapshot} combines two actual
calls, low/high decomposition, adjusted \(s_2\), and the frame in one Hoare
postcondition.
\theoremname{output_row_from_mode2_ntt_words} identifies
\identifier{KeygenM23ArithmeticSpec.output_row} exactly with the actual
full-invNTT/pointwise word representation.
The frozen paper ended at \identifier{STOP-KG-NTT}, lacking the odd-root
orthogonality/full-NTT convolution theorem and representation adapter needed
to transport that representation to native \(R_q\) negacyclic multiplication.

Post-freeze,
\theoremname{rq_mul_coeff_foldr_to_bigi},
\theoremname{full_ntt_montgomery_spectral_action},
\theoremname{full_ntt_montgomery_row_product},
\theoremname{keygen_mode2_row_product}, and
\theoremname{mode2_row_product_haetae_dot}
close this pure algebra and representation transport.
The two actual \identifier{avec} rows are also interpreted as the paper's
\(a\), and \(qj\) from the paper's \(j=(1,0)^T\) is added as a separate term
in the snapshot congruence.
The existing security-side \identifier{haetae_mode23_qj_vector}, however, is
a synthetic generator starting at coefficient 401, different from the paper's
\(qj[0,0]=64513\) and actual SHAKE128/rejection expansion.
The congruence above must therefore not simply be renamed to the existing
security object's \(A s=qj\pmod{2q}\), nor read as proving security-model
KG-1--KG-4. The current boundary is \CurrentGate; the frozen-paper boundary
remains unchanged
\cite{post-freeze-kg-rq-haetae-report,
post-freeze-kg-actual-avec-qj-report}.
\end{caution}

The latest standalone \sourcepath{KgFaithfulAugmentedModel.ec} does not
conceal this gap by bypassing it: it builds a separate faithful coefficient
carrier and instantiation with the actual matrix/secret.
\theoremname{actual_faithful_augmented_productE} expands the structured
coefficient product on an active row as the sum of a head, generated HAETAE
\identifier{poly_dot}, and adjusted error.
\theoremname{checked_mode2_parent_m23_finalize_faithful_augmented_paper_as_qj}
combines the actual sampler's paper \(a\), actual 3-term \(R_q\) row product,
and adjusted error into one expression, proving at every active coefficient
that \identifier{faithful_augmented_matrix_vector_product_coeff} is congruent
to \identifier{faithful_qj_coeff} mod-\(2q\).
This is \statusproved, with verdict \FaithfulAugmentedVerdict.
The product evaluator accepts only the actual width-6 matrix/secret and
consumes the head, middle three columns, and two identity columns.
Its mod-\(2q\) block-specific coefficient semantics avoid eliminating (qj)
in a mod-\(q\) ring. It is therefore neither the existing width-4 security
\identifier{matrix_vec_mul} theorem nor equivalence to synthetic security
ExpandVecA.

\subsection{The transcript where Sign and Verify meet}

In the pinned specification, the first observation Sign and Verify must
reconstruct in common is the message digest \(\mu\).
Restricting attention to the relevant portion,
\[
  \mu=H_{\mathrm{gen}}(\mathit{seedA},b_1,M),
\]
and the challenge input has the form
\[
  \rho=H\bigl(
      \operatorname{HighBits}_h(w),
      \operatorname{LSB}(\lfloor y_0\rceil),
      \mu\bigr)
\]
\cite{haetae-spec}. Sign reads \((\mathit{seedA},b_1)\) from the packed
secret key's 992-byte prefix; Verify reads the same bytes from public-key
memory.

The central equality currently proved at the raw/internal generated boundary is
\[
  \take_{32}\bigl(\mu_{\mathrm{Sign}}^{64}\bigr)
  =\mu_{\mathrm{Verify}}^{32}
\]
Absorbing these 32 bytes from the same challenge state at position~64 then
produces equal final states and positions. This is a relational theorem
directly including the actual generated hash/absorb procedures.

The entire \(\rho\) is not yet proved. In particular, actual Sign--Verify
agreement of \(\operatorname{HighBits}_h(w)\) and the LSB representation,
the complete challenge call, and distribution/rejection analysis remain
separate obligations. The closed \(\mu\) square therefore cannot be extended
to the complete challenge or to full \(\delta_{\mathrm{Sign}}=0\).

\subsection{Signature representation is a partial isomorphism on a canonical subset}

A mode~2 signature has 1474 bytes, with the actual extracted layout
\[
\underbrace{[0,32)}_{256\text{ challenge bits}}
\;\Vert\;
\underbrace{[32,1056)}_{1024\text{ signed low bytes}}
\;\Vert\;
\underbrace{[1056,1058)}_{\text{two size deltas}}
\;\Vert\;
\underbrace{[1058,1474)}_{\text{HBZ/h payload and padding}}
\]
Let \(C_{\mathrm{prefix}}\) be the set of canonical inputs in which every bit
word is \(0\) or \(1\) and every low word lies in the signed-byte range.
For the actual prefix packer and unpacker,
\[
  \Decode_{\mathrm{prefix}}\circ
  \Encode_{\mathrm{prefix}}
  =\operatorname{id}_{C_{\mathrm{prefix}}}
\]
is currently proved as partial correctness.

The domain restriction is substantive. On noncanonical words, encoding may
discard information or decoding may return a different representative.
In algebraic terms, \(C_{\mathrm{prefix}}\) is a chosen system of
representatives, and the round trip is a section--retraction relation on
that system.

This equality is not yet available for the entire signature space.
Size metadata, HBZ and \(h\) payloads, canonical zero padding, and success
and uniqueness of the complete parser remain open.

\subsection{Which algebraic facts belong to which layers}

\begin{center}
\begin{tabularx}{\textwidth}{L{31mm} Y L{30mm}}
\toprule
Object & Required proposition & Current status \\
\midrule
\(R_q\)-module KeyGen & Equality between Jasmin operations and the paper's
KeyGen distribution &
\statuspartial / top-level goal \\
actual KeyGen snapshot & Low/high decomposition of matrix/finalize results
and snapshot-only mod-\(2q\) congruence &
\statuspartial\ (KG-2-like boundary proved; NTT product bridge blocked) \\
packed key & \(\pi_{992}(sk)=\pi_{992}(vk)\) &
\statusproved\ (partial correctness) \\
raw \(\mu\) transcript & Equality of the first 32 bytes of Sign's 64-byte
output and Verify's 32 bytes &
\statusproved\ (generated raw/internal) \\
challenge suffix & Equal states after absorbing \(\mu_{32}\) from the same position~64 &
\statusproved\ (local zero-loss) \\
signature prefix & Decode-after-encode is the identity on canonical inputs &
\statusproved\ (prefix only) \\
HBZ/rANS suffix & Inverse relation of the actual full encoder and decoder &
\statusproved\ (production wrapper and fixed all-6 Hoare success; termination unproved) \\
complete signature scheme & KeyGen/Sign/Verify refinement and security composition &
\statusspecified / \statusblocked \\
\bottomrule
\end{tabularx}
\end{center}

The next section explains the logical distinctions between “partial correctness”,
“local”, “generated”, and “actual” in this table.
